% 定义底色和线条颜色
\definecolor{covercolor}{HTML}{315758}
\definecolor{linecolor}{HTML}{D6E19A}

% 书脊宽度（603页 × 0.05mm/页 ≈ 3cm）
\newlength{\spinewidth}
\setlength{\spinewidth}{3cm}

% 定义统一宽度
\newlength{\patternwidth}
\setlength{\patternwidth}{14.4cm}

\newfontfamily{\SenatorTall}{SenatorTall.otf}
\newfontfamily{\RebarBold}{rebar-bold-1.ttf}
\newfontfamily{\Trajan}{TrajanPro-Regular.otf}

\newcommand{\reversedpattern}{
    \begin{tikzpicture}[x=0.5\patternwidth/6, y=0.5\patternwidth/6,
        armfill/.style={fill=covercolor,draw=linecolor,line width=2.2pt},
        armline/.style={draw=linecolor,line width=2.2pt},
        arm/.pic={
            \draw[##1]
                (0,0) |- (4,4)
                -- ++(0,-\fpeval{4*(2-sqrt(2))})
                -- (45:4) -- cycle;
        }
    ]

        % ═══ 底层：光晕 ═══
        \draw[line width=3pt,linecolor,fill=covercolor] (-6,-6) rectangle (6,6);
        \foreach \i in {5.8, 5.75, ..., 3.9} {
            \def\tmp{\fpeval{100*(\i/5.8)}}
            \fill[covercolor!\tmp!white] circle[radius=\i];
        }

        % ═══ 8臂填充 ═══
        \foreach \angle in {0,45,...,315} {
            \pic[rotate=\angle] {arm={armfill}};
        }

        % ═══ 顶层：指数衰减高光 ═══
        \begin{scope}
            \clip (-6,-6) rectangle (6,6);
            \foreach \i in {0,1,...,100} {
                \pgfmathsetmacro{\radius}{\i/50*5.0}
                \pgfmathsetmacro{\opacity}{0.15 * exp(-\radius/1.2)}
                \fill[white, opacity=\opacity] (0,0) circle (\radius);
            }
        \end{scope}

        % ═══ 8臂边框 ═══
        \foreach \angle in {0,45,...,315} {
            \pic[rotate=\angle] {arm={armline}};
        }

    \end{tikzpicture}
}

% 封面部分
\pagecolor{covercolor}
\thispagestyle{empty}

\cleardoublepage
\thispagestyle{empty}

% 保存旧宽度，扩展：后封(21cm) + 书脊(3cm) + 前封(21cm) = 45cm
\newdimen\savepagewidth
\savepagewidth=\pdfpagewidth
\pdfpagewidth=\dimexpr 2\savepagewidth + \spinewidth\relax

% 偏移量 = 后封 + 书脊，minipage 推到这里才开始
\newdimen\coverx
\coverx=\dimexpr\savepagewidth + \spinewidth\relax

% 书脊
\begin{tikzpicture}[remember picture,overlay]
    \draw[linecolor,line width=0.8pt]
        ([xshift=\savepagewidth]current page.north west) --
        ([xshift=\savepagewidth]current page.south west);
    \draw[linecolor,line width=0.8pt]
        ([xshift=\coverx]current page.north west) --
        ([xshift=\coverx]current page.south west);
    \node[
        rotate=-90,
        anchor=center,
        text=linecolor,
        font=\RebarBold\fontsize{18}{22}\selectfont
    ] at ([xshift=\dimexpr\savepagewidth+0.5\spinewidth\relax]current page.west)
        {ZY6412 \hspace{1.2em} Topology \hspace{1.2em} 2nd};
\end{tikzpicture}

% 封面内容：用 \hspace* 推到前封区域
\hspace*{\coverx}%
\begin{minipage}{\textwidth}
\begingroup
    \centering

    \color{white}
    \vspace*{0.2cm}

    % TOPOLOGY标题
    {\centering
    \SenatorTall
    \fontsize{60}{0}\selectfont
    \noindent\makebox[\textwidth][s]{%
    \raisebox{1em}{\scalebox{3}[2.9]{T}}\hfill%
    \raisebox{1em}{\scalebox{3}[2.9]{O}}\hfill%
    \raisebox{1em}{\scalebox{3}[2.9]{P}}\hfill%
    \raisebox{1em}{\scalebox{3}[2.9]{O}}\hfill%
    \raisebox{1em}{\scalebox{3}[2.9]{L}}\hfill%
    \raisebox{1em}{\scalebox{3}[2.9]{O}}\hfill%
    \raisebox{1em}{\scalebox{3}[2.9]{G}}\hfill%
    \raisebox{1em}{\scalebox{3}[2.9]{Y}}%
    }
    \par}
    \vspace{-1cm}

    % SECOND EDITION
    \noindent
    \begin{minipage}{\patternwidth}
        \centering
        \RebarBold
        \color{linecolor}\fontsize{20}{0}\selectfont
        S\hspace{0.95em}E\hspace{0.95em}C\hspace{0.95em}O\hspace{0.95em}N\hspace{0.95em}D%
        \hfill%
        E\hspace{0.95em}D\hspace{0.95em}I\hspace{0.95em}T\hspace{0.95em}I\hspace{0.95em}O\hspace{0.95em}N%
    \end{minipage}
    \vspace{1.7cm}

    \reversedpattern

    \vspace{2em}
    \vfill

    % JAMES R. MUNKRES
    {\centering
    \Trajan
    \color{linecolor}\fontsize{38}{0}\selectfont
    \makebox[\patternwidth][s]{JAMES\hspace{0.3em} R.\hspace{0.3em} MUNKRES}
    \par}
    \vspace{0.1cm}

\endgroup
\end{minipage}

\clearpage
\pdfpagewidth=\savepagewidth
\nopagecolor
